% =====================================================================
%  preamble.tex
%
%  Packages, page geometry, theorem environments, and color setup.
%  Edit this file to change the document's overall look.
% =====================================================================

% --- Encoding & fonts -------------------------------------------------
\usepackage[T1]{fontenc}
\usepackage[utf8]{inputenc}
\usepackage{lmodern}
\usepackage{microtype}

% --- Page geometry ----------------------------------------------------
\usepackage[
  letterpaper,
  inner=1in, outer=1.4in,
  top=1in,   bottom=1in,
  marginparwidth=1.1in,
  marginparsep=4mm,
  headsep=14pt
]{geometry}

% --- Math -------------------------------------------------------------
\usepackage{amsmath}
\usepackage{amssymb}
\usepackage{amsthm}
\usepackage{mathtools}
\usepackage{bm}

% --- Graphics & color -------------------------------------------------
\usepackage{graphicx}
\graphicspath{{figures/}}
\usepackage[dvipsnames,table]{xcolor}
\usepackage{pgfplots}
\usepackage{tikz}
\usetikzlibrary{calc,arrows.meta,decorations.pathreplacing,patterns}

\pgfplotsset{compat=1.18}

% --- Tables, lists, captions -----------------------------------------
\usepackage{booktabs}
\usepackage{array}
\usepackage{enumitem}
\setlist{itemsep=2pt,topsep=2pt}
\usepackage{caption}
\captionsetup{font=small,labelfont=bf}

% --- Headers & footers ------------------------------------------------
\usepackage{fancyhdr}
\setlength{\headheight}{14pt}
\addtolength{\topmargin}{-2pt}
\pagestyle{fancy}
\fancyhf{}
\fancyhead[LE,RO]{\thepage}
\fancyhead[RE]{\nouppercase{\leftmark}}
\fancyhead[LO]{\nouppercase{\rightmark}}
\renewcommand{\headrulewidth}{0.4pt}

\fancypagestyle{plain}{%
  \fancyhf{}\fancyfoot[C]{\thepage}\renewcommand{\headrulewidth}{0pt}}

% --- Index ------------------------------------------------------------
\usepackage{imakeidx}

% --- Bibliography (BibTeX) -------------------------------------------
% Default: classic BibTeX (\bibliography{} in main.tex).
% To switch to biblatex+biber instead, comment out \bibliographystyle/
% \bibliography in main.tex and uncomment the two lines below, then
% replace \bibliography{bibliography} with \printbibliography.
% \usepackage[backend=biber,style=alphabetic]{biblatex}
% \addbibresource{bibliography.bib}

% --- Hyperref (load late) --------------------------------------------
\usepackage[
  colorlinks=true,
  linkcolor=NavyBlue,
  citecolor=Maroon,
  urlcolor=NavyBlue,
  bookmarksnumbered=true,
  bookmarksopen=true,
  pdfencoding=auto
]{hyperref}

% --- Theorem environments --------------------------------------------
% Numbered within chapters; shared counter for "theorem-like" results.
\theoremstyle{plain}
\newtheorem{theorem}{Theorem}[chapter]
\newtheorem{proposition}[theorem]{Proposition}
\newtheorem{lemma}[theorem]{Lemma}
\newtheorem{corollary}[theorem]{Corollary}

\theoremstyle{definition}
\newtheorem{definition}[theorem]{Definition}
\newtheorem{example}[theorem]{Example}
\newtheorem{exercise}{Exercise}[chapter]
\newtheorem{problem}[exercise]{Problem}

\theoremstyle{remark}
\newtheorem*{remark}{Remark}
\newtheorem*{note}{Note}

% A "counterexample" is highlighted because the book uses it
% systematically as a tag in the TOC.
\newtheoremstyle{counterex}%
  {\topsep}{\topsep}%               above/below
  {\itshape}%                       body font
  {}%                               indent
  {\bfseries\color{Maroon}}%        head font
  {.}%                              punctuation
  {.5em}%                           space after head
  {}%                               head spec
\theoremstyle{counterex}
\newtheorem{counterexample}[theorem]{Counterexample}

% Optional "warning" callout, useful in subsections marked Counterexample.
\newtheoremstyle{warningstyle}%
  {\topsep}{\topsep}{}{}%
  {\bfseries\color{BurntOrange}}{:}{.5em}{}
\theoremstyle{warningstyle}
\newtheorem*{warning}{Warning}

% --- Floats -----------------------------------------------------------
\usepackage{float}
\usepackage{wrapfig}

% --- Margin notes that stay put (used by \sectag) --------------------
\usepackage{marginnote}

% --- Cleveref (load after hyperref) ----------------------------------
\usepackage[capitalize,noabbrev]{cleveref}

% --- Misc -------------------------------------------------------------
\usepackage{epigraph}
\setlength{\epigraphwidth}{0.55\textwidth}
\renewcommand{\epigraphsize}{\small}
\renewcommand{\epigraphflush}{flushright}

% Tighter chapter title spacing (book class default is generous).
\usepackage{titlesec}
\titleformat{\chapter}[display]
  {\normalfont\sffamily\huge\bfseries}
  {\chaptertitlename\ \thechapter}{12pt}{\Huge}
\titlespacing*{\chapter}{0pt}{-20pt}{30pt}

\titleformat{\section}
  {\normalfont\sffamily\Large\bfseries}{\thesection}{1em}{}
\titleformat{\subsection}
  {\normalfont\sffamily\large\bfseries}{\thesubsection}{1em}{}

% Show paragraph and subparagraph numbers in TOC if the author wants
% deeper detail later (set tocdepth higher to expose them).
\setcounter{tocdepth}{2}     % part, chapter, section, subsection
\setcounter{secnumdepth}{3}  % number through subsection
